\section{Caminho de dados}
\label{sec:datapath}

A Figura~\ref{fig:datapath} mostra o caminho de dados da CPU. Três escolhas de
projeto organizam o desenho:

\begin{enumerate}[nosep]
  \item \textbf{A RAM só é acessada por \sig{MAR} e \sig{MBR}.} O endereço da
        memória é sempre o conteúdo de \sig{MAR}, e o dado de escrita é sempre o
        conteúdo de \sig{MBR}. Isso deixa os sinais da RAM estáveis durante um
        estado inteiro da UC, o que é necessário para a RAM síncrona.
  \item \textbf{A ALU tem entradas registradas.} Os operandos são copiados para
        \sig{ALU\_A} e \sig{ALU\_B} (com o comando \sig{CMD}) no estado de
        decodificação; no estado seguinte a ALU, que é combinacional, já
        apresenta resultado e flags estáveis, e a UC só precisa guardá-los.
  \item \textbf{\sig{MBR} centraliza os resultados.} Resultado da ALU, dado lido
        da memória, imediato de \instr{ldi} e valor de \instr{mov} passam todos
        por \sig{MBR} antes de ir para o registrador destino. Assim, a escrita
        no banco de registradores tem uma única origem e acontece sempre no
        mesmo estado (WRITE).
\end{enumerate}

\begin{figure}[H]
\centering
\newcommand{\muxat}[4]{% nome, x, y, meia-altura
  \draw[thick, fill=gray!12] (#2-0.22,#3-#4) -- (#2+0.22,#3-#4+0.2) -- (#2+0.22,#3+#4-0.2) -- (#2-0.22,#3+#4) -- cycle;
  \coordinate (#1-out) at (#2+0.22,#3);
  \coordinate (#1-in) at (#2-0.22,#3);}
\newcommand{\stub}[3]{% x_entrada, y, texto
  \draw[-{Latex[length=1.6mm]}, thick] (#1-0.55,#2) node[left, font=\tiny\ttfamily, inner sep=1pt] {#3} -- (#1,#2);}
\begin{tikzpicture}[scale=0.95, transform shape,
  reg/.style={draw, thick, fill=blue!6, minimum width=1.4cm, minimum height=0.7cm, font=\small\ttfamily},
  bus/.style={-{Latex[length=1.8mm]}, thick},
  lbl/.style={font=\tiny\ttfamily, fill=white, inner sep=0.8pt}]

  % ---------------- linha de cima: PC, SP, IR, UC ----------------
  \muxat{mpc}{1.2}{4.0}{0.75}
  \stub{0.98}{4.45}{PC+1}
  \stub{0.98}{4.15}{PC+2}
  \stub{0.98}{3.85}{MBR}
  \stub{0.98}{3.55}{REG(rx)}
  \node[reg] (pc) at (2.7,4.0) {PC};
  \draw[bus] (mpc-out) -- (pc.west);
  \node[lbl, right=2pt of pc.east, align=left] {$\to$ MAR\\$\to$ PC+1, PC+2};

  \node[reg] (sp) at (7.0,4.0) {SP};
  \stub{6.3}{4.0}{SP$\pm$1}
  \node[lbl, below=1pt of sp] {$\to$ MAR (SP, SP+1)};

  \node[reg] (ir) at (10.2,4.0) {IR};
  \node[draw, thick, fill=yellow!12, rounded corners=2pt, minimum width=2.0cm, minimum height=0.7cm,
        font=\small, align=center] (uc) at (13.0,4.0) {UC\\[-2pt]\scriptsize decode + FSM};
  \draw[bus] (ir.east) -- (uc.west);

  % ---------------- banco de registradores ----------------
  \node[draw, thick, fill=blue!6, minimum width=1.4cm, minimum height=3.0cm, align=center, font=\small\ttfamily]
       (rf) at (0.8,0.6) {Ra\\[2pt]Rb\\[2pt]Rc\\[2pt]Rd};
  \node[font=\scriptsize, above=0pt of rf] {registradores};

  % ---------------- ALU ----------------
  \node[reg] (aa) at (3.4,1.4) {ALU\_A};
  \muxat{mb}{2.45}{-0.2}{0.5}
  \node[reg] (ab) at (3.4,-0.2) {ALU\_B};
  \draw[bus] ([yshift=0.8cm]rf.east) -- node[lbl, above] {REG(rx)} (aa.west);
  \draw[bus] ([yshift=-0.55cm]rf.east) -- node[lbl, above, pos=0.4] {REG(ry)} ([yshift=0.25cm]mb-in);
  \draw[-{Latex[length=1.6mm]}, thick] (2.45,-1.35) node[below, font=\tiny\ttfamily, inner sep=1pt] {1, 0} -- (2.45,-0.6);
  \draw[bus] (mb-out) -- (ab.west);

  \draw[thick, fill=orange!15] (5.0,2.0) -- (6.4,1.25) -- (6.4,-0.05) -- (5.0,-0.8) -- (5.0,0.25) -- (5.3,0.6) -- (5.0,0.95) -- cycle;
  \node[font=\small\bfseries] at (5.85,0.6) {ALU};
  \node[font=\tiny\ttfamily, anchor=north] at (5.7,-0.45) {CMD};
  \draw[bus] (aa.east) -- (5.0,1.4);
  \draw[bus] (ab.east) -- (5.0,-0.2);

  \node[reg, fill=red!8] (fl) at (5.7,-2.0) {FLAGS};
  \draw[bus] (5.7,-0.4) -- (fl.north);
  \draw[bus] (fl.west) -- ++(-0.6,0) node[left, lbl] {LEDs, UC (saltos)};

  % ---------------- MBR ----------------
  \muxat{mmbr}{8.7}{1.8}{0.85}
  \draw[bus] (6.4,0.6) -- (7.6,0.6) |- node[lbl, pos=0.75, above] {ALU\_R} (8.48,1.15);
  \stub{8.48}{2.45}{DOUT}
  \stub{8.48}{2.0}{REG(rx)}
  \stub{8.48}{1.6}{REG(ry)}
  \node[reg] (mbr) at (10.2,1.8) {MBR};
  \draw[bus] (mmbr-out) -- (mbr.west);

  % ---------------- MAR ----------------
  \muxat{mmar}{8.7}{-1.3}{1.0}
  \stub{8.48}{-0.6}{PC+1, próx. PC}
  \stub{8.48}{-1.05}{DOUT}
  \stub{8.48}{-1.5}{REG(ry)}
  \stub{8.48}{-1.95}{SP, SP+1}
  \node[reg] (mar) at (10.2,-1.3) {MAR};
  \draw[bus] (mmar-out) -- (mar.west);

  % ---------------- RAM ----------------
  \node[draw, thick, fill=green!8, minimum width=1.9cm, minimum height=3.6cm, align=center]
       (ram) at (13.0,0.4) {\textbf{RAM}\\\scriptsize $256\times8$\\\scriptsize síncrona};
  \draw[bus] (mbr.east) -- node[lbl, above] {DIN} (mbr.east -| ram.west);
  \draw[bus] (mar.east) -- node[lbl, above] {ADDR} (mar.east -| ram.west);
  \draw[bus] (14.5,-0.4) node[right, lbl] {WE} -- (ram.east |- 0,-0.4);

  % DOUT -> IR (e para os muxes de MBR e MAR, indicados como entradas rotuladas)
  \draw[bus] (ram.east |- 0,1.6) -- (14.5,1.6) node[right, lbl, align=left] {DOUT\\$\to$ MBR, MAR} -- (14.5,3.1) -- (10.2,3.1) -- (ir.south);
  \node[lbl] at (12.0,3.1) {DOUT};

  % write-back MBR -> registradores
  \draw[bus] (mbr.south) -- (10.2,0.9) -- (11.3,0.9) -- (11.3,-2.9) -- (0.8,-2.9) node[lbl, pos=0.5, above] {MBR $\to$ Rx (estado WRITE)} -- (rf.south);
\end{tikzpicture}
\caption[Caminho de dados da CPU]{Caminho de dados da CPU\protect\footnotemark}
\label{fig:datapath}
\end{figure}
\footnotetext{Os sinais de controle gerados pela UC foram omitidos para
clareza; os multiplexadores são selecionados pela instrução decodificada e pelo
estado.}

O que cada registrador recebe em cada estado, para cada instrução, está
detalhado na Tabela~\ref{tab:microops}, na seção da máquina de estados.
Os únicos somadores fora da ALU são os incrementos de \sig{PC} (\sig{PC+1},
\sig{PC+2}) e de \sig{SP} (\sig{SP+1}, \sig{SP-1}), exatamente as exceções
permitidas pelo enunciado. Todas as operações das instruções --- inclusive
\instr{inc} e \instr{dec} --- são feitas pela ALU.
